% Options for packages loaded elsewhere
\PassOptionsToPackage{unicode}{hyperref}
\PassOptionsToPackage{hyphens}{url}
%
\documentclass[
]{article}
\usepackage{amsmath,amssymb}
\usepackage{iftex}
\ifPDFTeX
  \usepackage[T1]{fontenc}
  \usepackage[utf8]{inputenc}
  \usepackage{textcomp} % provide euro and other symbols
\else % if luatex or xetex
  \usepackage{unicode-math} % this also loads fontspec
  \defaultfontfeatures{Scale=MatchLowercase}
  \defaultfontfeatures[\rmfamily]{Ligatures=TeX,Scale=1}
\fi
\usepackage{lmodern}
\ifPDFTeX\else
  % xetex/luatex font selection
\fi
% Use upquote if available, for straight quotes in verbatim environments
\IfFileExists{upquote.sty}{\usepackage{upquote}}{}
\IfFileExists{microtype.sty}{% use microtype if available
  \usepackage[]{microtype}
  \UseMicrotypeSet[protrusion]{basicmath} % disable protrusion for tt fonts
}{}
\makeatletter
\@ifundefined{KOMAClassName}{% if non-KOMA class
  \IfFileExists{parskip.sty}{%
    \usepackage{parskip}
  }{% else
    \setlength{\parindent}{0pt}
    \setlength{\parskip}{6pt plus 2pt minus 1pt}}
}{% if KOMA class
  \KOMAoptions{parskip=half}}
\makeatother
\usepackage{xcolor}
\usepackage{color}
\usepackage{fancyvrb}
\newcommand{\VerbBar}{|}
\newcommand{\VERB}{\Verb[commandchars=\\\{\}]}
\DefineVerbatimEnvironment{Highlighting}{Verbatim}{commandchars=\\\{\}}
% Add ',fontsize=\small' for more characters per line
\newenvironment{Shaded}{}{}
\newcommand{\AlertTok}[1]{\textcolor[rgb]{1.00,0.00,0.00}{\textbf{#1}}}
\newcommand{\AnnotationTok}[1]{\textcolor[rgb]{0.38,0.63,0.69}{\textbf{\textit{#1}}}}
\newcommand{\AttributeTok}[1]{\textcolor[rgb]{0.49,0.56,0.16}{#1}}
\newcommand{\BaseNTok}[1]{\textcolor[rgb]{0.25,0.63,0.44}{#1}}
\newcommand{\BuiltInTok}[1]{\textcolor[rgb]{0.00,0.50,0.00}{#1}}
\newcommand{\CharTok}[1]{\textcolor[rgb]{0.25,0.44,0.63}{#1}}
\newcommand{\CommentTok}[1]{\textcolor[rgb]{0.38,0.63,0.69}{\textit{#1}}}
\newcommand{\CommentVarTok}[1]{\textcolor[rgb]{0.38,0.63,0.69}{\textbf{\textit{#1}}}}
\newcommand{\ConstantTok}[1]{\textcolor[rgb]{0.53,0.00,0.00}{#1}}
\newcommand{\ControlFlowTok}[1]{\textcolor[rgb]{0.00,0.44,0.13}{\textbf{#1}}}
\newcommand{\DataTypeTok}[1]{\textcolor[rgb]{0.56,0.13,0.00}{#1}}
\newcommand{\DecValTok}[1]{\textcolor[rgb]{0.25,0.63,0.44}{#1}}
\newcommand{\DocumentationTok}[1]{\textcolor[rgb]{0.73,0.13,0.13}{\textit{#1}}}
\newcommand{\ErrorTok}[1]{\textcolor[rgb]{1.00,0.00,0.00}{\textbf{#1}}}
\newcommand{\ExtensionTok}[1]{#1}
\newcommand{\FloatTok}[1]{\textcolor[rgb]{0.25,0.63,0.44}{#1}}
\newcommand{\FunctionTok}[1]{\textcolor[rgb]{0.02,0.16,0.49}{#1}}
\newcommand{\ImportTok}[1]{\textcolor[rgb]{0.00,0.50,0.00}{\textbf{#1}}}
\newcommand{\InformationTok}[1]{\textcolor[rgb]{0.38,0.63,0.69}{\textbf{\textit{#1}}}}
\newcommand{\KeywordTok}[1]{\textcolor[rgb]{0.00,0.44,0.13}{\textbf{#1}}}
\newcommand{\NormalTok}[1]{#1}
\newcommand{\OperatorTok}[1]{\textcolor[rgb]{0.40,0.40,0.40}{#1}}
\newcommand{\OtherTok}[1]{\textcolor[rgb]{0.00,0.44,0.13}{#1}}
\newcommand{\PreprocessorTok}[1]{\textcolor[rgb]{0.74,0.48,0.00}{#1}}
\newcommand{\RegionMarkerTok}[1]{#1}
\newcommand{\SpecialCharTok}[1]{\textcolor[rgb]{0.25,0.44,0.63}{#1}}
\newcommand{\SpecialStringTok}[1]{\textcolor[rgb]{0.73,0.40,0.53}{#1}}
\newcommand{\StringTok}[1]{\textcolor[rgb]{0.25,0.44,0.63}{#1}}
\newcommand{\VariableTok}[1]{\textcolor[rgb]{0.10,0.09,0.49}{#1}}
\newcommand{\VerbatimStringTok}[1]{\textcolor[rgb]{0.25,0.44,0.63}{#1}}
\newcommand{\WarningTok}[1]{\textcolor[rgb]{0.38,0.63,0.69}{\textbf{\textit{#1}}}}
\setlength{\emergencystretch}{3em} % prevent overfull lines
\providecommand{\tightlist}{%
  \setlength{\itemsep}{0pt}\setlength{\parskip}{0pt}}
\setcounter{secnumdepth}{-\maxdimen} % remove section numbering
\ifLuaTeX
  \usepackage{selnolig}  % disable illegal ligatures
\fi
\usepackage{bookmark}
\IfFileExists{xurl.sty}{\usepackage{xurl}}{} % add URL line breaks if available
\urlstyle{same}
\hypersetup{
  hidelinks,
  pdfcreator={LaTeX via pandoc}}

\author{}
\date{}

\begin{document}

\section{Panduan Penggunaan Repo Proposal
LaTeX}\label{panduan-penggunaan-repo-proposal-latex}

Repo ini dipakai untuk menulis, merevisi, dan meninjau proposal skripsi
LaTeX secara kolaboratif menggunakan GitHub.

\subsection{Tujuan Workflow}\label{tujuan-workflow}

Alur kerja repo ini dibuat agar: - file skripsi diedit di laptop
masing-masing - hasil perubahan masuk ke GitHub dengan rapi - setiap
revisi diperiksa lewat Pull Request sebelum digabung ke branch utama -
proses compile dilakukan secara lokal, bukan di Overleaf

\subsection{Struktur File Utama}\label{struktur-file-utama}

\begin{itemize}
\tightlist
\item
  \texttt{main.tex} - file utama untuk compile dokumen
\item
  \texttt{frontmatter.tex} - bagian awal dokumen
\item
  \texttt{chapter1.tex} - Bab 1
\item
  \texttt{chapter2.tex} - Bab 2
\item
  \texttt{chapter3.tex} - Bab 3
\item
  \texttt{appendices.tex} - lampiran
\item
  \texttt{references.bib} - daftar pustaka
\item
  \texttt{figures/} - folder gambar
\end{itemize}

\subsection{Cara Compile Lokal}\label{cara-compile-lokal}

Repo ini menggunakan \texttt{XeLaTeX} dan \texttt{Biber}.

Jalankan perintah berikut di folder project:

\begin{Shaded}
\begin{Highlighting}[]
\ExtensionTok{xelatex}\NormalTok{ main.tex}
\ExtensionTok{biber}\NormalTok{ main}
\ExtensionTok{xelatex}\NormalTok{ main.tex}
\ExtensionTok{xelatex}\NormalTok{ main.tex}
\end{Highlighting}
\end{Shaded}

Jika memakai VS Code, bisa gunakan extension LaTeX Workshop. Jika
memakai TeXstudio, pastikan engine yang dipilih adalah \texttt{XeLaTeX}
dan bibliografi menggunakan \texttt{Biber}.

\subsection{Alur Jika Ratih Sudah Mengubah File di
Laptop}\label{alur-jika-ratih-sudah-mengubah-file-di-laptop}

Kalau Ratih sudah mengedit file di laptopnya dan ingin memasukkan
perubahan ke repo GitHub, langkahnya seperti ini.

\subsubsection{1. Masuk ke folder
project}\label{masuk-ke-folder-project}

\begin{Shaded}
\begin{Highlighting}[]
\BuiltInTok{cd}\NormalTok{ Proposal\_LaTeX}
\end{Highlighting}
\end{Shaded}

\subsubsection{2. Ambil versi terbaru dari
repo}\label{ambil-versi-terbaru-dari-repo}

Lakukan ini dulu sebelum mulai push, supaya tidak tertinggal perubahan
terbaru.

\begin{Shaded}
\begin{Highlighting}[]
\FunctionTok{git}\NormalTok{ pull origin main}
\end{Highlighting}
\end{Shaded}

\subsubsection{3. Buat branch baru}\label{buat-branch-baru}

Jangan langsung kerja di branch \texttt{main}.

Contoh:

\begin{Shaded}
\begin{Highlighting}[]
\FunctionTok{git}\NormalTok{ checkout }\AttributeTok{{-}b}\NormalTok{ revisi{-}bab2{-}ratih}
\end{Highlighting}
\end{Shaded}

\subsubsection{4. Edit file yang
diperlukan}\label{edit-file-yang-diperlukan}

Contoh file yang diedit: - \texttt{chapter1.tex} - \texttt{chapter2.tex}
- \texttt{chapter3.tex} - \texttt{references.bib}

\subsubsection{5. Compile lokal dulu}\label{compile-lokal-dulu}

Sebelum push ke GitHub, pastikan dokumen masih bisa dicompile.

\begin{Shaded}
\begin{Highlighting}[]
\ExtensionTok{xelatex}\NormalTok{ main.tex}
\ExtensionTok{biber}\NormalTok{ main}
\ExtensionTok{xelatex}\NormalTok{ main.tex}
\ExtensionTok{xelatex}\NormalTok{ main.tex}
\end{Highlighting}
\end{Shaded}

Kalau masih error, selesaikan dulu error lokalnya sebelum lanjut.

\subsubsection{6. Cek file yang berubah}\label{cek-file-yang-berubah}

\begin{Shaded}
\begin{Highlighting}[]
\FunctionTok{git}\NormalTok{ status}
\end{Highlighting}
\end{Shaded}

\subsubsection{7. Simpan perubahan ke
Git}\label{simpan-perubahan-ke-git}

\begin{Shaded}
\begin{Highlighting}[]
\FunctionTok{git}\NormalTok{ add .}
\FunctionTok{git}\NormalTok{ commit }\AttributeTok{{-}m} \StringTok{"revisi bab 2 oleh Ratih"}
\end{Highlighting}
\end{Shaded}

\subsubsection{8. Push branch ke GitHub}\label{push-branch-ke-github}

\begin{Shaded}
\begin{Highlighting}[]
\FunctionTok{git}\NormalTok{ push origin revisi{-}bab2{-}ratih}
\end{Highlighting}
\end{Shaded}

\subsubsection{9. Buat Pull Request}\label{buat-pull-request}

Setelah branch berhasil di-push: - buka repo GitHub - akan muncul tombol
untuk membuat Pull Request - pilih base branch \texttt{main} - isi judul
dan deskripsi perubahan - submit Pull Request

\subsection{Alur Saat Kamu Menerima Pull
Request}\label{alur-saat-kamu-menerima-pull-request}

Kalau Ratih atau kolaborator lain sudah membuat Pull Request, langkah
review yang disarankan:

\subsubsection{1. Baca perubahan di
GitHub}\label{baca-perubahan-di-github}

Periksa file apa saja yang berubah dan apakah revisinya memang sesuai
kebutuhan.

\subsubsection{2. Tarik branch itu ke laptop jika perlu dicek
langsung}\label{tarik-branch-itu-ke-laptop-jika-perlu-dicek-langsung}

\begin{Shaded}
\begin{Highlighting}[]
\FunctionTok{git}\NormalTok{ fetch origin}
\FunctionTok{git}\NormalTok{ checkout revisi{-}bab2{-}ratih}
\end{Highlighting}
\end{Shaded}

\subsubsection{3. Compile lokal}\label{compile-lokal}

\begin{Shaded}
\begin{Highlighting}[]
\ExtensionTok{xelatex}\NormalTok{ main.tex}
\ExtensionTok{biber}\NormalTok{ main}
\ExtensionTok{xelatex}\NormalTok{ main.tex}
\ExtensionTok{xelatex}\NormalTok{ main.tex}
\end{Highlighting}
\end{Shaded}

\subsubsection{4. Review isi revisi}\label{review-isi-revisi}

Cek hal berikut: - apakah isi tulisan sudah benar - apakah format LaTeX
tetap rapi - apakah sitasi dan daftar pustaka tetap normal - apakah
tidak ada file penting yang terhapus

\subsubsection{5. Jika sudah aman, merge Pull
Request}\label{jika-sudah-aman-merge-pull-request}

Merge dilakukan ke \texttt{main}.

\subsection{Aturan Kerja yang
Disarankan}\label{aturan-kerja-yang-disarankan}

\begin{itemize}
\tightlist
\item
  Jangan edit langsung di branch \texttt{main}.
\item
  Satu branch untuk satu jenis revisi.
\item
  Sebisa mungkin satu orang fokus ke satu bab agar konflik kecil.
\item
  Jika mengubah \texttt{references.bib}, beri tahu anggota lain karena
  file ini rawan konflik.
\item
  Selalu compile lokal sebelum commit dan sebelum merge.
\item
  Gunakan pesan commit yang jelas.
\end{itemize}

\subsection{Contoh Nama Branch}\label{contoh-nama-branch}

\begin{itemize}
\tightlist
\item
  \texttt{revisi-bab1-ratih}
\item
  \texttt{tambah-referensi-bab2}
\item
  \texttt{perbaiki-abstrak}
\item
  \texttt{rapikan-frontmatter}
\end{itemize}

\subsection{Contoh Pesan Commit}\label{contoh-pesan-commit}

\begin{itemize}
\tightlist
\item
  \texttt{revisi\ latar\ belakang\ bab\ 1}
\item
  \texttt{menambah\ sitasi\ pada\ bab\ 2}
\item
  \texttt{perbaiki\ format\ abstrak}
\item
  \texttt{rapikan\ daftar\ pustaka}
\end{itemize}

\subsection{Jika Repo Belum Pernah Diunggah ke
GitHub}\label{jika-repo-belum-pernah-diunggah-ke-github}

Kalau repo ini belum dihubungkan ke GitHub, jalankan sekali saja:

\begin{Shaded}
\begin{Highlighting}[]
\FunctionTok{git}\NormalTok{ init}
\FunctionTok{git}\NormalTok{ add .}
\FunctionTok{git}\NormalTok{ commit }\AttributeTok{{-}m} \StringTok{"initial commit proposal latex"}
\FunctionTok{git}\NormalTok{ branch }\AttributeTok{{-}M}\NormalTok{ main}
\FunctionTok{git}\NormalTok{ remote add origin }\OperatorTok{\textless{}}\NormalTok{url{-}repo{-}github}\OperatorTok{\textgreater{}}
\FunctionTok{git}\NormalTok{ push }\AttributeTok{{-}u}\NormalTok{ origin main}
\end{Highlighting}
\end{Shaded}

Setelah itu, workflow berikutnya tinggal pakai branch dan Pull Request
seperti langkah di atas.

\end{document}
